\documentclass{article}
\usepackage[a4paper, total={6in,8in}]{geometry}
\usepackage{outlines, enumerate}
\usepackage{amsmath,amssymb,amsthm}




\begin{document}
	\section{Brute Force Classification}
	\subsection{Supervised Machine Learning}
	\begin{outline}
		\1 Input space $X$, Output space $Y$
		\1 Training data: given many instances $\mathcal{D} = \{(\vec{x}_i, y_i)\}^N_{i=1}$
		\1 Learn a function $X \rightarrow Y$ which is expressed by parameters
		\1 Test data: Given new $\vec{x}_{N+1}$, use function to predict $y_{N+1}$
		\1 Evaluation: Check whether prediction $y_{N+1}$ matches the ground truth $y^*_{N+1}$
		\1 $y^*_{N+1}$ is not known to the black box (held-out)
	\end{outline}
	\subsection{Text Categorization}
	\begin{outline}
		\1 Given a text document, assign it to one of $K$ categories.
		\1 Consider how you're deciding the category and how you would explain it to the computer.
		\1 Input space: documents, Output: one of $K$ categories
		\1 Training data: documents and their category (e.g. sports, news)
		\1 Learn a categorization function
		\1 Test data: Given a new document, use function to predict category
		\1 Evaluation: Check whether predicted category is correct
		\1 Evaluation measure: how many correct predictions?
	\end{outline}
	\subsubsection{Binary Text Categorization Task}
	\begin{outline}
		\1 There is only one category.
		\1 Predict whether the document is in the category or not (pos/neg)
		\1 Same: Predict the set of documents that are in this category.
		\1 Input space: documents, Output: $\{0, 1\}$ (1=positive)
		\1 Training data: documents and whether they are in this category (e.g. sports or not)
		\1 Learn a categorization function
		\1 Evaluation measure: how many correct predictions?
	\end{outline}
	\subsubsection{One-vs-Rest Classification}
	\begin{outline}
		\1 Input space: documents, Output: one of $K$ categories
		\1 Redefine the task to learn $K$ binary categorization functions
		\1 Learn $K$ categorization functions (one per category)
		\1 Can a document be positive for multiple categories? If yes, which category to choose?
		\2 Each categorization function bases the decision on a \textbf{Classification Score} or \textbf{Confidence}
		\2 Assign the document to the category with the highest confidence
	\end{outline}
	\subsubsection{Supervised Document Classification}
	\begin{outline}
		\1 Represent documents with feature vectors
		\1 Example content features:
			\2 $f_1{d}:$ 1 if doc $d$ contains the word "ball"; 0 otherwise
			\2 $f_2(d):$ 1 if doc $d$ contains the world "goal"; 0 otherwise
			\2 $f_3(d):$ 1 if doc $d$ contains the word "election"; 0 otherwise
		\1 Change the parameters in the function to get more predictions right
		\1 Iteratively update parameters (Training loop)
		\1 Then fix the parameters, make one prediction in the test set
		\1 Sometimes: very good on the train set, but errors on the test set
		\1 {\bf Note: Never train on the test data!}
			\2 Never test on the training data (unless you call it training error)
			\2 Never leak the test data into your training process
	\end{outline}
	\subsection{A Simple Categorization Model}
	\subsubsection{Linear Hyperplane Model}
	\begin{outline}
		\1 Feature vectors in space $X$ (number of features = dimension of space)
		\1 Find a hyperplane that separates positive from negative examples
		\1 To classify a new example, see if it is on the "positive" side of the hyperplane
			\2 Weight vector $\vec{w}$ points towards positive side
		\1 {\it The hyperplane has one fewer dimension than the vector space}
	\end{outline}
	\subsubsection{Hyperplane Parameters}
	\begin{outline}
		\1 Hyperplane is represented by vector $\vec{w}$
		\1 $\vec{w}$ is orthogonal to the hyperplane
		\1 $\vec{w}$ has the same dimensionality as $\vec{x}_i$
		\1 If $\vec{w} \cdot \vec{x}_i > 0$ then example $i$ is positive
		\1 Inner product $\vec{w} \cdot \vec{x}_i = \sum_jw_jx_{ij}$ over feature dimensions $j$.
	\end{outline}
	\subsubsection{Inner Products and Cosines}
	\begin{outline}
		\1 If $\vec{w} \cdot \vec{x} > 0$ then example $i$ is positive
		\1 Same as: $\angle(\vec{w}, \vec{x}_i) < 90\deg$
		\1 Same as: $\cos(\vec{w}, \vec{x}_i) > 0$ \hspace{3em}$\cos(\vec{a}, \vec{b}) = \frac{\vec{a} \cdot \vec{b}}{||\vec{a}||\:||\vec{b}||}$
	\end{outline}
	\subsubsection{Hard vs Soft Classification}
	\begin{outline}
		\1 Sometimes we want to tolerate classification errors in training data to get a more robust hyperplane
		\1 We don't want outliers and errors to dictate our model.
		\1 \textbf{Hard classification} is strict; the hyperplane exactly and entirely divides the space into positive and negative
		\1 \textbf{Soft classification} allows for some errors, typically to allow for a simpler hyperplane. This also prevents outliers and errors from causing issues with classification
	\end{outline}
	\subsection{The Intercept of the Model}
	\subsubsection{Changing the Threshold with "The Intercept"}
	\begin{outline}
		\1 Changing the threshold $b$: $\vec{w} \cdot \vec{x}_1 > b$ \hspace{2em} The intercept $b$
	\end{outline}
	\subsubsection{"The Intercept" Trick}
	\begin{outline}
		\1 It is easy to forget to add the intercept for $\vec{w} \cdot \vec{x}_i > b$.
		\1 Rewrite as $\vec{w} \cdot \vec{x}_i + (-b)\cdot 1 > 0$, recognize the inner products
	\end{outline}
\end{document}